\begin{corrige}{4-11}

\begin{enumerate}

\item $f$ est bien une fonction de densit\'e: $f\geq0$ car la fonction indicatrice d'un ensemble est positive, et 
$$
\int_{-\infty}^{+\infty} f(s) ds =
\int_{-\infty}^{+\infty} \frac1{b-a} \mathds{1}_{[a,b]} (s) ds =
\frac1{b-a} \int_a^b ds = \frac{b-a}{b-a} = 1.
$$

\item Soit $Z$ une variable al\'eatoire qui suit la loi uniforme $\mathcal{U}(0,1)$ vue en cours. On v\'erifie que:
$$ 
X = a + (b-a) Z.
$$
On a alors, en utilisant la relation ci-dessus et la d\'efinition des fonctions de r\'epartition:
$$
F_X (t) = F_Z \left ( \frac{t-a}{b-a} \right ),
$$
et on en d\'eduit que:
$$
F_X (t) = \frac{t-a}{b-a}, \quad a\leq t \leq b,
$$
avec $F_X(t)=0$ si $t<a$ et $F_X(t)=1$ si $t>b$.

\item On a:
$$
P ( c \leq X \leq d ) = F_X(d) - F_X(c) = \frac{d-a}{b-a} - \frac{c-a}{b-a} = \frac{d - c}{b-a}.
$$
On peut dire que la loi est uniforme car la probabilit\'e que $X$ soit dans un intervalle donn\'e ne d\'epend que de la longueur de cet intervalle, et non de sa localisation sur la droite r\'eelle.

\item On peut faire le calcul direct en utilisant la d\'efinition du cours. Toutefois, il est plus rapide d'utiliser la relation $X = a + (b-a) Z$. En effet:
$$
E[X] = E [ a + (b-a) Z ] = E [ a ] + (b-a) E [ Z ] = a + \frac{b-a}2 = \frac{a+b}2.
$$

\item De m\^eme, en utilisant $a + (b-a) Z$:
$$
V(X) = V ( a + (b-a) Z ) = (b-a)^2 \, V(Z) = \frac{(b-a)^2}{12}.
$$

\end{enumerate}

\end{corrige}
